% ================================================================
%  CHAPTER 1 -- INTRODUCTION
%  Status: DRAFT v1 (2026-05-18).
% ================================================================
\chapter{Introduction}
\label{ch:introduction}

% ----------------------------------------------------------------
\section{Motivation}
\label{sec:motivation}

Bayesian inference provides a principled framework for reasoning about
unknown parameters under uncertainty. Given observed data and a
statistical model, it combines prior beliefs with the likelihood to
produce a posterior distribution that encodes our knowledge about
the parameters after observing the data. However, for all but a small class of so-called conjugate models the posterior is analytically intractable, and
inference must proceed by computational approximation.

The dominant approximation is Markov chain Monte Carlo (MCMC), reviewed in Chapter~\ref{ch:mcmc}, which
constructs a Markov chain whose stationary distribution is the
posterior and draws samples by simulating it~\citep{brooks2011handbook, robert2004monte}. MCMC is reliable and
theoretically well understood: under broad conditions, the chain converges to the exact
posterior, and the approximation error is reduced simply by running it
for longer. It carries, however, a structural cost. An MCMC chain
targets the posterior for one specific dataset. When a new dataset is
observed the entire procedure must typically be repeated from scratch, because
the chain transfers no information from one inference problem to the
next. In settings where the same model must be fitted to hundreds or
thousands of distinct datasets, as arises routinely in spatial
statistics and the analysis of replicated
experiments~\citep{sainsburydale2024nbe, sainsburydale2025graph}, the total cost scales linearly in the number of
problems, with no economy of scale.

\emph{Amortised} inference takes the opposite trade-off~\citep{zammitmangion2025neural}. Rather than
solving each inference problem independently, one invests
computational effort once to learn a reusable mapping from
data to posterior summaries. Inference for any new dataset then costs
only a single evaluation of that mapping. A \emph{neural Bayes
estimator} (NBE)~\citep{sainsburydale2024nbe} realises this idea by training
a neural network to approximate the Bayes estimator, the function
that maps data to the posterior summary minimising a chosen loss
function (developed in Chapter~\ref{ch:amortization}). Neural networks are typically used here because they are good function approximators, have reliable black-box training methods, and are fast to run in forward mode. The
network is trained on data simulated from the model, so the procedure
is likelihood-free: it requires only the ability to simulate, not to
evaluate the likelihood. Once trained, the network can quickly generate parameter estimates (such as posterior means) by applying what is effectively just linear algebra to a new dataset -- this is known as a forward pass of the network.

The training of an NBE has at its heart a Monte Carlo
computation. The network is fitted by minimising an empirical average
of the loss over simulated parameter--data pairs, and this empirical
average is a Monte Carlo estimate of the population Bayes risk. Its
gradient, the quantity that drives training, is likewise a Monte
Carlo estimate of the population gradient, and is therefore a random
quantity with a variance.

This thesis takes that observation as its starting point. If training
an NBE involves Monte Carlo integration, then the
techniques developed over decades to reduce the variance of Monte
Carlo estimators can be applied not to the inference itself, but to the
training of the estimator. A training gradient with lower
variance is a better-behaved optimisation signal, and a better-behaved
signal should produce a better estimator for the same computational
budget.

The primary technique developed here is \emph{Rao-Blackwellisation} (Chapter~\ref{ch:rao-blackwell}). The Rao-Blackwell theorem states that conditioning an estimator on
additional information cannot increase its variance under a convex
loss. Applied to the training gradient of an NBE, it
implies that if part of the parameter vector has a tractable
conditional posterior given the rest, that part can be integrated
out of the loss analytically rather than sampled. The resulting
training gradient is provably no noisier than the standard one, and is
typically much less noisy. The conditional structure this exploits,
conditional conjugacy, is the same structure that makes Gibbs
sampling possible, and it is common in hierarchical models.

Another variance-reduction technique, \emph{importance sampling}, is also investigated in Chapter~\ref{ch:importance-sampling}. Unlike Rao-Blackwellisation, importance sampling does not require a tractable conditional posterior; it only requires simulated training examples to be drawn from a proposal distribution, then reweighted back such that the network still learns the unbiased estimator. This makes it more broadly applicable in principle, but also less automatic. Its benefit depends on the proposal placing simulation effort where the loss, or the training-gradient noise, is actually concentrated. In the autoregressive experiment this condition fails. The estimator is already close to the exact posterior across most of the prior, and the proposal emphasises regions that are not the dominant source of its remaining error. Reweighting therefore does not improve accuracy; when pushed further, it lowers the effective sample size and adds noise.



% ----------------------------------------------------------------
\section{Related work}
\label{sec:related-work}

\subsection{Simulation-based inference}
\label{sec:rw-sbi}

The technique of producing estimators here belongs to the broader category of \emph{simulation-based
inference}: the family of methods for fitting models that can be
simulated from but whose likelihood is intractable or too expensive to
evaluate~\citep{cranmer2020frontier}. The classical representative of this method is
approximate Bayesian computation, which compares simulated and
observed data through summary statistics and a tolerance and underlies
much of the applied likelihood-free literature. Over the past decade
the field has been reshaped by neural networks, which replace
hand-crafted summaries and accept--reject schemes with learned
mappings between data and inferential targets~\citep{zammitmangion2025neural}.
The neural approaches are not uniformly superior. A systematic
benchmark finds no single best algorithm, and shows that the choice of
performance metric and the sample efficiency of the training scheme
matter as much as the architecture~\citep{lueckmann2021benchmarking};
but they have made amortised, reusable inference practical for models
that were previously out of reach.

\subsection{Normalising flows}
\label{sec:rw-flows}

Neural simulation-based methods are often distinguished by what they
learn. \emph{Neural posterior estimation} learns the full posterior
density, \emph{neural likelihood estimation} learns the
likelihood, and \emph{neural ratio
estimation} learns a likelihood ratio for model comparison~\citep{papamakarios2019snl}. Each
typically represents the learned distribution with a \emph{normalising
flow}, a bijective neural transformation of a simple base density whose
tractable Jacobian yields an exact likelihood~\citep{papamakarios2017maf,
papamakarios2021flows}. The \textsc{BayesFlow} framework couples such
an invertible network with a summary network that learns informative
statistics directly from simulations, giving a globally amortised
posterior sampler~\citep{radev2020bayesflow, radev2023bayesflow}. These
methods return a complete posterior, at the cost of training and
storing a generative model of it.

\subsection{Neural Bayes estimation and amortised inference}
\label{sec:rw-nbe}

Neural Bayes estimation, the setting of this thesis, is the most
economical point on this spectrum: it learns only a point summary of
the posterior, the Bayes estimator under a chosen loss, trading
the full distribution of the density-estimation methods for speed and
simplicity. The estimator in the form used here is due to
\citet{sainsburydale2024nbe}, who establish the decision-theoretic
framework, the permutation-invariant architecture for replicated data,
and applications to weakly identified and highly parameterised models;
the approach has since been extended to irregular spatial data with
graph neural networks~\citep{sainsburydale2025graph}. It sits within
the broader programme of amortised inference surveyed by
\citet{zammitmangion2025neural}: methods that pay a one-off training
cost up front and thereafter return estimates for new datasets in single fast forward passes, in a fraction of a second where likelihood-based inference could
take hours. That survey spans point estimation, approximate Bayesian
inference, learned summary statistics, and likelihood approximation;
but, like the works above, it treats the variance of the training
objective as incidental rather than as a quantity to be controlled.
That variance is the gap this thesis takes up.

\subsection{The simulation cost of likelihood-free inference}
\label{sec:rw-efficiency}

A separate strand of work attacks the efficiency of likelihood-free
inference from the simulation side. Because these methods lean on
enormous numbers of model simulations, much can be saved by spending
the simulation budget wisely. Multifidelity and multilevel schemes
combine cheap, approximate simulators with the expensive model of
interest and allocate computation across fidelities to minimise cost
at a fixed accuracy~\citep{warne2022multilevel}, an allocation that can
be made adaptive, learning the relationship between fidelities on the
fly to approach the optimal trade-off~\citep{prescott2024multifidelity}.
This line is complementary to the present work: it reduces the number
and cost of the simulations that feed inference, whereas this
thesis reduces the variance of the training gradient extracted
from a given set of simulations; the two could in principle be
combined.

\subsection{Variance reduction for stochastic gradients}
\label{sec:rw-variance}

The technique developed here belongs to the literature on variance
reduction for stochastic optimisation, which includes control variates
and related estimators~\citep{wang2013variance}. 
Rao-Blackwellisation
in particular has been used to reduce the variance of gradient
estimators elsewhere in machine learning, notably for the
straight-through Gumbel-softmax estimator~\citep{paulus2021raoblackwell}
and for stochastic gradients of discrete
distributions~\citep{liu2019raoblackwell}.
Those methods Rao-Blackwellise the gradient with respect to discrete latent
variables or reparameterisation noise; the contribution of this thesis
differs in what is conditioned on. The training gradient of a neural
Bayes estimator is Rao-Blackwellised with respect to a conjugate
block of the model parameters, exploiting the conditional conjugacy of
the inference model itself: the same structure that makes Gibbs
sampling possible.

% ----------------------------------------------------------------
\section{Contributions}
\label{sec:contributions}

The contributions of this thesis are as follows.

\begin{enumerate}
  \item \textbf{Frames the variance of the training gradient as an
  explicit, reducible target.} A NBE is trained by
  minimising an empirical (Monte Carlo) Bayes
  risk~\citep{sainsburydale2024nbe}, so its training gradient is itself
  a Monte Carlo estimator with a variance: a property of how the
  loss is sampled, not of the inference problem. Recognising this
  brings the Monte Carlo variance-reduction toolbox to bear on the
  training of these estimators, which their literature has not
  previously addressed (Chapter~\ref{ch:amortization}).

  \item \textbf{Derives the \emph{Rao-Blackwellised loss}:} When a block of
  the parameter vector has a tractable conditional posterior, that
  block is integrated out of the loss analytically. For squared-error
  loss the construction reduces to substituting the conditional
  posterior mean for the sampled target, and the residual conditional
  variance contributes no gradient (Chapter~\ref{ch:rao-blackwell}).

  \item \textbf{Proves that the resulting gradient-variance reduction is
  positive semi-definite} (Rao-Black\-wellisation can never increase the
  variance of the training gradient, for any model or architecture), and shows that for a linear network the same reduction carries to
  the fitted network weights themselves (Chapter~\ref{ch:rao-blackwell}).

  \item \textbf{Demonstrates the method empirically on Bayesian linear
  regression.} Across a grid of problem dimensions and batch sizes, the
  Rao-Blackwellised estimator attains a lower loss than the standard NBE in every configuration, closing about half of the gap to
  the theoretical minimum, the Bayes-optimal floor computed by Gibbs sampling, and matching the
  standard NBE at a quarter of the minibatch size
  (Chapter~\ref{ch:rao-blackwell}).

  \item \textbf{Contrasts Rao-Blackwellisation with importance sampling,} a
  second variance-reduction technique, on an autoregressive model,
  where importance sampling returns a negative result: establishing
  that its benefit is contingent on the problem, where the
  Rao-Blackwell guarantee is unconditional
  (Chapter~\ref{ch:importance-sampling}).
\end{enumerate}

% ----------------------------------------------------------------
\section{Thesis outline}
\label{sec:outline}

Chapter~\ref{ch:mcmc} reviews Markov chain Monte Carlo, the
per-dataset incumbent against which amortised inference is contrasted,
and collects three running case studies --- the normal,
linear-regression, and autoregressive models --- used later as Bayes-optimal
references for neural estimators. Chapter~\ref{ch:amortization} develops the Bayes estimator
framework, defines neural Bayes estimators, and establishes the
central observation that their training uses Monte Carlo integration.
Chapter~\ref{ch:rao-blackwell} is the core of the thesis: it derives
the Rao-Blackwellised loss, proves the gradient-variance reduction,
analyses its effect on the converged loss, and tests the method on
Bayesian linear regression. Chapter~\ref{ch:importance-sampling}
develops importance sampling as a contrasting technique and applies it
to an autoregressive model. Chapter~\ref{ch:conclusion} summarises the
findings and discusses limitations and directions for further work.
